% ============================================================
%  Remark: why training from scratch outperforms fine-tuning
%  under the K=V constraint
% ============================================================

% ---------- Notation (consistent with main paper) -----------
%  M     = W_O (W_V^pre - W_K^pre)       (cross-talk matrix)
%  rho   = Condition C coverage (0--1)
%  theta_snap = pretrained model with W_V <- W_K, rest fixed
% ------------------------------------------------------------

\begin{remark}[Fine-tuning vs.\ Training from Scratch]
\label{rem:ft-vs-scratch}
Kayyam et al.\ establish empirically that K$=$V tying succeeds when
the model is trained from scratch: the tied model matches untied
quality, confirming the constraint imposes no fundamental
expressivity barrier.
However, \emph{retrofitting} the K$=$V constraint onto a pretrained
model via fine-tuning is a qualitatively harder problem, for a
reason made precise by Theorem~2.

By Theorem~2 Part~(ii), the number of steps to reach
$\mathbb{E}[\mathcal{L}(\theta_T)] - \mathcal{L}^*_{\mathrm{tied}} \leq \varepsilon$
is
\begin{equation}
  T \;=\; O\!\left(
    \frac{L(\theta_0) - \mathcal{L}^*_{\mathrm{tied}}}{\mu\,\eta\,\varepsilon}
  \right),
  \label{eq:T-bound}
\end{equation}
so convergence time is directly proportional to the initial loss gap.
The two training regimes differ only in $\mathcal{L}(\theta_0)$:

\begin{itemize}
  \item \textbf{Training from scratch.}
        $\theta_0$ is a random initialisation;
        $\mathcal{L}(\theta_0) \approx \log|V|$, independent of $\rho$.

  \item \textbf{Fine-tuning.}
        $\theta_0 = \theta_{\mathrm{snap}}$, the pretrained model with
        $W_V \leftarrow W_K$.  This snap operation removes the
        value-specific subspace that $W_V^{\mathrm{pre}}$ encodes,
        incurring an immediate loss spike.  Experiments~2 and~3 measure
        this gap directly:

        \begin{center}
        \begin{tabular}{lrrrr}
          \toprule
          Model & $\rho$ & $1-\rho^2$ &
            PPL (untied) & PPL (snap) \\
          \midrule
          OPT-125M & 0.904 & 0.183 & 59.2  & 7{,}569  \\
          OPT-350M & 0.677 & 0.542 & 47.7  & 35{,}116 \\
          \bottomrule
        \end{tabular}
        \end{center}

        The snap PPL grows with $1-\rho^2$: the less the pretrained
        model satisfies Condition~C, the more the snap damages it.
\end{itemize}

Substituting into \eqref{eq:T-bound}, the step-count ratio is
\begin{equation}
  \frac{T_{\mathrm{FT}}}{T_{\mathrm{SC}}}
  \;\approx\;
  \frac{\mathcal{L}(\theta_{\mathrm{snap}}) - \mathcal{L}^*_{\mathrm{tied}}}
       {\log|V| - \mathcal{L}^*_{\mathrm{tied}}},
  \label{eq:ratio}
\end{equation}
which is directly computable from the measured snap PPL and the
model's random-init cross-entropy.
For OPT-350M the ratio exceeds $100{\times}$, consistent with the
near-zero PPL recovery observed after 20{,}000 fine-tuning steps
in Experiment~3.

The conclusion is that fine-tuning under K$=$V is not a viable
strategy for models with low $\rho$: the snap penalty swamps any
reasonable training budget.
$\rho$ is therefore a \emph{retrofit feasibility diagnostic} ---
it predicts, before any fine-tuning begins, whether the
K$=$V constraint can be absorbed within a given step budget.
\end{remark}
